\section{Experimental Methodology}
\label{sec:experiments}

This section describes the implemented M1--M1.6 development protocol. It is
intended to make the pilot auditable and to freeze the comparison design before
the proposed \rcqasyn{} components are evaluated. The three datasets are
development datasets, not the complete final benchmark.

\subsection{Datasets, Registry, and Splits}

Dataset definitions are versioned in \texttt{configs/datasets/}. Source and
acquisition records live under \texttt{data/manifests/}, while each run persists the
resolved TRAIN profile and split manifest in its \texttt{dataset/} directory. The
current development set covers moderate, severe, and extreme binary class
imbalance and includes one mixed-type dataset. Exact characteristics and the
committed splits appear in \cref{tab:development-datasets}.

\begin{table}[t]
  \centering
  \caption{Development datasets and committed deterministic stratified splits.
  Class counts are full-dataset negative/positive counts; IR is
  majority/minority. These are pilot datasets, not the complete final
  benchmark.}
  \label{tab:development-datasets}
  \small
  \begin{tabularx}{\linewidth}{@{}lrrrrrX@{}}
    \toprule
    Dataset & Rows & Num. & Cat. & Neg./Pos. & IR & TRAIN / TEST \\
    \midrule
    page-blocks0 & 5,472 & 10 & 0 & 4,913 / 559 & 8.79 & 4,377 / 1,095 \\
    Ecoli4       &   336 &  7 & 0 &   316 / 20  & 15.80 & 269 / 67 \\
    Abalone19    & 4,174 &  7 & 1 & 4,142 / 32  & 129.44 & 3,340 / 834 \\
    \bottomrule
  \end{tabularx}
\end{table}


Each dataset is split once by deterministic stratified sampling under root
seed 42. Profiles, observed numeric summaries, categorical domains, class
proportions, and reference examples are computed from TRAIN only. REAL TEST is
reserved for downstream evaluation: it is not used in prompt construction,
deterministic validation, judge prompts, model fitting, or threshold design.
The run manifests record \texttt{no\_test\_leakage=true}. The current final-benchmark
plan is broader but remains unevaluated, as summarized in
\cref{tab:planned-benchmark}.

\begin{table}[t]
  \centering
  \caption{Planned final-benchmark roles. Inclusion and protocol remain to be
  frozen; listing a dataset does not imply that a final experiment exists.}
  \label{tab:planned-benchmark}
  \small
  \begin{tabularx}{\linewidth}{@{}lXl@{}}
    \toprule
    Dataset & Intended stress-test role & Status \\
    \midrule
    Default Credit Card & Mild imbalance / control & Planned \\
    HTRU2 & Moderate real-world imbalance & Planned \\
    Ecoli4 & Small-data scarcity & Development only \\
    Musk & High dimensionality & Planned \\
    Mammography & Severe imbalance & Planned \\
    Bank Marketing & Mixed numerical/categorical features & Planned \\
    Abalone19 & Extreme imbalance & Development only \\
    \bottomrule
  \end{tabularx}
  \par\smallskip
  \footnotesize APS and Satimage-2 remain optional candidates in the research
  plan; neither is claimed as part of a frozen benchmark.
\end{table}


\subsection{P1 Generation Protocol}
\label{sec:p1-protocol}

P1 uses the local OpenAI-compatible Ollama endpoint convention
\texttt{http://127.0.0.1:11434/v1}; this deployment choice is not a scientific
contribution. The generator is \texttt{gemma4:e4b}. For each dataset, the configuration
requests 200 records in explicit class-conditioned batches: 100 negative and
100 positive, with batch size five. The TRAIN context contains a compact
profile and four deterministic reference records, two per class. Generation
uses prompt version v2, temperature 0.7, a 2,048-token output limit, and a
120-second timeout.

Prompt v2 requests one JSON object with a single \texttt{records} array. This
transport wrapper was introduced to match the provider's JSON-object response
format; it is separate from the feature schema and does not enforce field
validity. Raw provider output is stored before parsing. Parse errors, wrong
outer shapes, and item-level failures are retained rather than clipped,
imputed, or silently regenerated.

The root seed and configured provider base seed are both 42. A cryptographic
derivation based on SHA-256 maps
\begin{equation}
  (\text{root seed},\text{dataset},\text{requested class},\text{batch index},
  \text{base seed})
  \longmapsto \text{provider seed},
  \label{eq:seed-derivation}
\end{equation}
with the result reduced to the provider's unsigned 32-bit range. The complete
seed plan is computed before generation and checked for collisions across all
scheduled dataset/class/batch tuples. Provider traces bind every candidate to
the corresponding batch index and derived seed.

\subsection{Paired Baselines}
\label{sec:baselines}

All three baselines operate on the same persisted parseable candidate pool.
They differ only in the acceptance policy.

\paragraph{B1---raw generation.}
B1 accepts every parseable candidate. It performs no filtering, repair, or
regeneration; parse failures remain outside the candidate pool and are
reported separately.

\paragraph{B2---deterministic filtering.}
B2 runs five independent checks: exact declared schema and primitive numeric
types; TRAIN-observed categorical domains; the declared missing-value policy;
agreement between the generated target and requested class; and observed TRAIN
numeric ranges. Schema, invalid-domain, rejected-missing, and target-condition
failures have HIGH severity. Observed-range findings have MEDIUM severity and
are warnings by default. The configured policy rejects HIGH and CRITICAL
evidence. In addition, B2 rejects every repeated canonical candidate after its
first occurrence. It does not learn a fusion rule and does not modify records.

\paragraph{B3---single LLM judge.}
B3 makes one judge call per parseable candidate using \texttt{gemma4:e4b}, the same
model used for generation. The judge runs at temperature 0 with prompt version
v1, strict JSON-object parsing, a static threshold of 0.7, and rejection of
malformed outputs. A candidate is accepted only when the parsed verdict is
PASS or WARN and its score reaches the threshold. The score is self-reported
and uncalibrated; B3 uses no ensemble, evidence fusion, external factual
verification, repair, or adaptive threshold. Sharing the generator model
creates a possible correlated-error and self-evaluation limitation rather than
a benefit.

\subsection{Intrinsic and Downstream Metrics}

Numeric features are compared between REAL TRAIN and accepted synthetic data
with the two-sample Kolmogorov--Smirnov statistic and one-dimensional
Wasserstein distance. Categorical features, including the target, use total
variation distance and Jensen--Shannon divergence. Pearson and Spearman
dependence error is the mean absolute difference between corresponding
off-diagonal correlations over numerical feature pairs.

Diversity outputs include the number of unique canonical rows, exact duplicate
count and rate, base-two categorical entropy, numerical variance relative to
REAL TRAIN, and the number of unique categorical combinations. Numeric,
categorical, and duplicate views are also reported by target class. Empty and
degenerate cases carry explicit statuses and null values rather than synthetic
zeros.

Downstream utility uses train-real/test-real (TRTR) as a reference and
train-synthetic/test-real (TSTR) as the synthetic-data condition. The
implemented classifiers are logistic regression with standardized numerical
features and a \texttt{liblinear} solver, and a 50-tree random forest; categorical
features are imputed and one-hot encoded. Metrics are F1, average precision
(reported under both \texttt{average\_precision} and \texttt{pr\_auc}), Matthews correlation
coefficient, balanced accuracy, and AUROC. No single metric is sufficient:
distributional similarity does not imply useful decision boundaries,
downstream utility does not reveal redundancy or constraint failures, and
record validity does not establish dataset coverage.

\subsection{Copying and Support Diagnostics}
\label{sec:diagnostics}

M1.5 computes three independent exact-copy indicators for every raw candidate:
whether it repeats an earlier synthetic record, equals a TRAIN reference shown
in its prompt, or equals a TRAIN record not shown in the prompt. These
categories distinguish generation-scale redundancy from visible-example
anchoring and contact with unshown TRAIN data. Exact TRAIN matching is a
copying diagnostic; it is not called privacy leakage in this study.

Support diagnostics measure Euclidean nearest-neighbour distance to REAL TRAIN
after numerical z-score standardization with TRAIN mean and population
standard deviation. Zero-variance numeric columns are removed. Categorical
features are one-hot encoded over TRAIN domains with unit weight. The reports
include nearest TRAIN, nearest same-class TRAIN, and nearest other-class TRAIN
distance where defined. These distances are representation-dependent,
learning-free diagnostics with no validity threshold, probability
interpretation, or effect on B1/B2/B3 decisions.

\subsection{Blinded Human Evaluation}
\label{sec:human-evaluation}

The master audit contains 120 unique-prioritized records, 40 from each
development dataset. Its manifest records overlapping selection-stratum hits:
18 accepted by all three baselines, eight B2/B3 disagreements, seven B2
rejections, one B3 rejection, 18 high-distance and 16 low-distance records, 15
majority and 15 minority records, 15 non-copies, and seven reference copies.
Because these are overlapping priority hits, they are not intended to sum to
120.

Reviewer A receives all 120 records. Reviewer B receives a deterministic
40-record overlap selected with seed 42: 14 Abalone19, 14 Ecoli4, and 12
page-blocks0 records, balanced as 20 negative and 20 positive requested
classes. Reviewers see only an anonymous annotation ID, record ID, dataset,
requested class, serialized candidate, and the five human-entry fields. They
do not see B2/B3 decisions, B3 scores or reasons, copying flags, support
distances, provider metadata, or automated verdicts.

For blinded annotation, a human-valid record conforms to known schema and
domain semantics and is a plausible member of the requested class,
independently of redundancy within the synthetic dataset. Redundancy alone
does not make a record invalid. Reviewers assign a boolean validity label and
one or more controlled record-level failure types; severity, affected fields,
and \texttt{human\_notes} are optional. Agreement will be summarized with raw
agreement and Cohen's kappa on explicitly labelled overlap. Disagreements are
manually adjudicated and are never resolved automatically. Human annotation is
in progress, so no human-validity result appears in this draft.

\subsection{Reproducibility and Development Confounds}
\label{sec:reproducibility}

P1 was executed with Python 3.12.11. Each run stores resolved configuration,
dataset and split fingerprints, TRAIN profile, prompt hash and version,
generation plan, seed plan, provider traces, raw candidates, baseline
decisions/evidence, metrics, and diagnostics. The three P1 manifests identify
source commit \texttt{ab492bf294b563e4b165755692ac8edb8f3dd2be} but also record
\texttt{git\_dirty=true}; the persisted artifacts, not a claim of a clean historical
checkout, are therefore the primary record of the runs.

An earlier P0-v2 protocol reused one provider seed across same-class batches.
The provider replayed raw batches, creating an artificial 75\% exact duplicate
rate on each development dataset. The problem was diagnosed as a seed
confound, and the per-dataset/class/batch derivation in
\cref{eq:seed-derivation} was implemented with a pre-run uniqueness guard.
P0-v3 and P1 use the corrected protocol. Confounded P0-v2 diversity values are
not evidence about generator quality.

One further artifact discrepancy remains open: the Ecoli4 P1 run contains 201
parseable records although the configuration requested 200. The extra record
came from the persisted generation output and is not silently discarded in
the reports. Consequently, Ecoli4 rates in this draft use the recorded pool
denominator where stated, and the count anomaly is treated as a development
limitation requiring resolution before the final benchmark.
